\scanpage{76}
\deleted{(finally surjectivity is established) that $T_\alpha,T_\beta$ are in a $=$ connected component of $X$, there exists a ``chain'' of \uline{indices}
\[
i_0=i,\ i_1,\ldots,i_n=j
\]
such that for two consecutive indices $k=i_{\unsure{r|\nu}}$, $k'=i_{\unsure{r|\nu}+1}$, one has $T_k\subset T_{k'}$ or $T_{k'}\subset T_k$.}
\margin{Suppose moreover that the interiors of the $T_\alpha$ cover $X$ \uline{or} that the $T_\alpha$ are closed (\deleted{and} $A$ finite) Then $\Pi_0(A)\xrightarrow{\sim}\Pi_0(X)$. Better ---}

Suppose moreover \inserted{that the interiors of the $T_i$ cover $X$,} \deleted{that the $T_i$ s'} \inserted{\uline{This being so,}} I say that knowledge of the (for $u$ fixed) of the $g_{y,x}^\alpha(u)\in\pi_\alpha$ \inserted{($\alpha\in A,y,x\in T_\alpha$)} allows one to recover all the invariants $g_l(u)$ (or $f_l(u)$). Indeed, it suffices to see that the arrows $l_{y,x}^\alpha=p_\alpha(l_{y,x})$ ($\alpha\in I$, $x,y\in T_\alpha$, $p_\alpha:T_\alpha\hookrightarrow X$ the inclusion) \deleted{generate} form a system of generators of the groupoid $\Pi_1(X_K)$\deleted{\unsure{$x_0$}} --- \inserted{all} \unsure{the|some} \deleted{objects} vertices are obtained in any case, and one finds a bijective map on the $\Pi_0$ (of the groupoid generated by the generators $l_{y,x}^\alpha$ \inserted{($\alpha\in I,x,y\in T_{\alpha K}$) and the relations $l_{y,x}^{\beta}=\varphi_{\beta\alpha}(l_{y,x}^{\alpha})$} and of $\Pi_1(X_K)$), one must see that one has a surjective map on the $\Pi_1$.

When the interiors of the $T_{\inserted{\alpha}\deleted{\unsure{\alpha|i}}}$ cover $X$, it suffices to note that knowledge of the $g_{y,x}^{\alpha}(u)$ gives us the $g_{y,x}(u)$ (without superscript $\alpha$ !) for $x,y$ points ``nearby'' on $X$,
\margin{75}
